\section{各任务定义}
\label{appendix-definition}

\subsubsection{检测与跟踪}

作为自动驾驶系统的组成部分，检测与跟踪使自车能够定位和感知周围环境中的障碍物。具体而言，检测模块需要感知自车周围的障碍物，并提供障碍物的位置、尺寸及相对于自车的朝向等信息。此外，在跟踪模块中（检测模块的时序扩展），每个障碍物被分配一个唯一ID，该ID应在各帧之间保持一致。

\subsubsection{在线建图}

旨在替代离线标注的高精地图，在线建图负责建模驾驶环境的结构和语义信息，能够显著节省人力成本并扩展自动驾驶系统的适用范围。具体而言，将地图视为线段（Line Segment）的稀疏表示，每个道路元素对应一组稀疏线段，并带有特定类别，如分隔线、人行横道、道路边界和车道线。

\subsubsection{运动预测}

运动预测模块负责预测所有交通参与者可能的未来状态，以辅助自车的决策过程。通常，历史信息对于预测周围智能体的未来轨迹是必要的。此外，地图也是运动预测模块输入的重要组成部分。需要指出的是，聚合和提取周围环境的信息不仅有利于周围智能体的运动预测，也有利于自车的规划。

\subsubsection{规划}

作为自动驾驶流程接近末端的关键组件，规划旨在提供一条同时满足安全性和舒适性要求并符合运动学规律的行驶轨迹。在上述目标的指导下，规划模块将来自多个上游任务的所有信息（包括当前和可能的周围环境及智能体的未来状态）作为输入，并充分考虑不同类型的约束，以生成无碰撞、舒适且可控的轨迹。
